\chapter{The Nervous System}
\section{Get Rolling}
\textbf{We know that all the functions of a computer are controlled by the CPU. The monitor, the keyboard and the mouse all are connected to the CPU. Similarly, there is a control centre in the human body, which controls almost all the bodily functions. The brain acts as the control centre in the human body.}
\par
\vspace{0.1in}
We use our brain to do almost all our activities. Name four such activities. (One is given as an example.)
\begin{enumerate}
    \item Reading
    \item Thinking
    \item Doing
    \item Solving
    \item Playing
\end{enumerate}
Have you wondered how the brain controls functions of all the organs? The brain is attached to all the organs through nerves. The brain together with the nerves forms a well-organised system that controls all the bodily functions. This system is known as the \textbf{nervous system.}
\section{What to Know}
The nervous system and it's parts, reflex actions, and structure and functions of different sense organs
\section{\textbf{Learning Objectives}}
\begin{itemize}
    \item to describe the functions of the nervous system
    \item to list the different parts of he nervous system
    \item to describe the different parts of the brain
    \item to explain reflex actions
    \item to state the functions of the different sense organs
\end{itemize}
\section{I Explore}
\section{Nervous System}
The \textbf{nervous system} controls and coordinates all the functions of the body. We can see, hear, smell, think and learn because of the nervous system.
\vspace{0.1in}
\par
The nervous system has three main parts: the brain, and the spinal cord and nerves.
\subsection{Brain}
The brain, inside the skull, is the body's command centre. A fluid-filled space between the skull and the brain cushions it from injures and jerks. The brains various parts have specific tasks.
\vspace{0.1in}
\par
The \textbf{cerebrum}, the largest part, handles thinking, problem-solving, memory senses and body movements.
\vspace{0.1in}
\par
The \textbf{cerebellum} manages balance, coordination and movement during activities like biking or sports.
\vspace{0.1in}
\par
The \textbf{brainstem (medulla)} connects the brain to the body, controlling vital functions like breathing, heartbeat, blood pressure and sleep.
\vspace{1.4in}
\begin{enumerate}
    \item Use different coloured playdough for different brain parts.
    \item Shape cerebrum as the largest part, place cerebellum towards the back and create a thin, vertical structure for the brainstem.
    \item Attach labels to corresponding clay parts using glue or tape.
    \item Discuss functions and importance of each part of the brain after assembling the model.
\end{enumerate}
\vspace{1.6in}
\subsection{Spinal Cord}
The \textbf{spinal cord}, a long tube running through the spine, connects the brainstem to the lower back. It acts as a messenger, transmitting signals between the brain and the body. For example, when we wish the move our limbs, the brain sends signals down the spinal cord to instruct the muscles.
\subsection{Nerves}
\textbf{Nerves} are like bundles of long, thin threads in the body that carry signals between different body parts and the brain. They are composed of nerve cells called \textbf{neurons}, the basic building blocks of the nervous system.
\vspace{0.1in}
\par
There are three main types of nerves in the body.
\par
\textbf{Sensory nerves:} These nerves transmit information from our sense organs to the brain. For instance, when we touch something hot, the sensory nerves in the skin send a message to the brain, signalling that we have touched something hot and prompting us to move our hand away.
\vspace{0.1in}
\par
\textbf{Motor nerves:} These nerves convey signals from the brain to muscles, enabling movement and control. When we want to kick a ball, the brain sends signals via motor nerves to our leg muscles, instructing them to contract and kick the ball.
\vspace{0.1in}
\par
\textbf{Mixed nerves:} Most nerves in the body are mixed nerves, which means they carry both sensory and motor signals. They carry messages in two directions, sensory information from the body to the brain and motor commands from the brain to the muscles.
\vspace{0.3in}
\par
I wonder which type of nerves will help me regain control when I stumble!
\section{Reflex Actions}
\textbf{Reflex actions} are automatic, involuntary responses that do not involve thinking. For example, when we touch something hot, sensory nerves send a quick message to the spinal cord. The spinal cord instantly sends a message back through motor nerves, directing the muscles to react. The brain does not play any role here. This quick, automatic process acts like our body's safety system, protecting us from harm in dangerous situations. In this automatic process, there is a stimulus and a response. Here, the stimulus is touching a hot object, and the response is removing our hand instantly.
\vspace{2.0in}
\section{Sense Organs}
There are five sense organs: eyes, ears, nose, skin and tongue. These organs help us know about and feel the world around us.
\subsection{Eyes}
Eyes help us see by working together with different parts.
\vspace{0.1in}
\par
The coloured-part \textbf{iris} controls light by adjusting the pupil's size and the \textbf{pupil} lets light in. The dome-shaped and transparent \textbf{cornea} protects the eye's front.
\vspace{0.1in}
\par
The \textbf{lens} bends light to the back and the \textbf{retina} is like a special screen at the eye's back.
\vspace{0.1in}
\par
The \textbf{optic nerve} acts like a cable, sending messages from the eyes to the brain for understanding what we see.
\subsection{Ears}
We have a pair of ears that help us hear. The ears also help in maintaining the balance in the body. Let us learn about the different parts of the ear and the functions they perform.
\vspace{0.1in}
\par
The ear has three main parts: outer ear, middle ear and inner ear.
\vspace{0.1in}
\par
The \textbf{outer ear}, consisting of the pinna and ear canal, collects and directs sound waves.
\vspace{0.1in}
\par
The \textbf{middle ear}, with the eardrum and three ear bones, passes sound vibrations from the eardrum to the inner ear.
\vspace{0.1in}
\par
The \textbf{inner ear} has a cochlea that converts sound vibrations into signals. The brain understands these signals and enables hearing.
\vspace{2.8in}
\subsection{Nose}
How do you smell your favourite food? Through your nose! The nose, with its fine hair and muscus, traps dust particles in the air we breathe. The air enters through nostrils, leading to the nasal cavity with smell-sensitive cells. These cells detect smells and send messages to the brain.
\subsection{Tongue}
The tongue senses sweet, salty, bitter and sour tastes through taste buds all over its surface. These taste buds send messages to the brain.
\vspace{0.1in}
\par
Taste and smell signals combined create our perception of flavour.
\vspace{1.1in}
\subsection{Skin}
The skin is the largest organ and forms the outer covering of the body. It helps us feel, touch and sense heat, cold, pressure and pain. Under the surface of the skin, there are nerves that detect touch and send messages to the brain.
\section{Now I Know}
\begin{itemize}
    \item The \textbf{nervous system} controls and coordinates all the functions of the body.
    \item The brain helps us think, learn and remember things.
    \item The brain has three parts: cereburm, cerebellum and brainstem.
    \item The spinal cord runs through the centre of the spine, from the brainstem to the lower back. It helps our brain communicate with the rest of our body.
    \item \textbf{Nerves} are specialised structures in the human body that transmit signals or messages between different parts of the body and the brain. They are made up of nerve cells called \textbf{neurons.}
    \item There are three types of nerves: sensory nerves, motor nerves and mixed nerves.
    \item \textbf{Reflex actions} are quick and automatic responses that do not involve thinking. They are controlled by the spinal cord.
\end{itemize}